\section{回顾：图像分类与线性分类器}

本节课开始之前，讲者首先回顾了前一节课的核心内容。图像分类是计算机视觉中的基础任务：给定一张图片，将其映射到一组预定义标签中的某一个。

\begin{figure}[H]
\centering
\includegraphics[width=0.95\textwidth]{slides-images/slide-001.jpg}
\caption{图像分类回顾：将图像映射到标签集合中的正确类别\protect\footnotemark}
\end{figure}
\footnotetext{Slides 第2页。}

\subsection{图像分类的核心挑战}

图像分类面临多个困难：

\begin{itemize}
    \item \textbf{语义鸿沟}（Semantic Gap）：人类看到的是``猫''，计算机看到的是一个三维像素数组（$32 \times 32 \times 3 = 3072$ 个数值）
    \item \textbf{光照变化}（Illumination）：不同光照条件下像素强度差异巨大
    \item \textbf{形变}（Deformation）：猫等柔性物体可以扭曲弯折成各种形状
    \item \textbf{遮挡}（Occlusion）：物体可能被部分遮挡
    \item \textbf{背景混淆}（Background Clutter）：物体可能与背景融为一体
    \item \textbf{类内差异}（Intra-class Variation）：同一类别的不同实例外观差异很大
\end{itemize}

\begin{figure}[H]
\centering
\includegraphics[width=0.95\textwidth]{slides-images/slide-002.jpg}
\caption{图像分类面临的挑战：语义鸿沟、光照、形变、遮挡、背景混淆和类内差异\protect\footnotemark}
\end{figure}
\footnotetext{Slides 第3页。}

\subsection{数据驱动方法与线性分类器}

由于无法用简单的 if-else 规则处理所有情况，我们转向数据驱动方法。线性分类器的核心思想是：将图像展平为向量 $\mathbf{x} \in \mathbb{R}^{3072}$，通过权重矩阵 $W \in \mathbb{R}^{10 \times 3072}$ 进行矩阵乘法得到各类别的分数：

$$f(\mathbf{x}, W) = W\mathbf{x} + \mathbf{b}$$

\begin{itemize}
    \item $W$：权重矩阵，每一行代表一个类别的``模板''
    \item $\mathbf{b}$：偏置向量，每个类别一个偏置值
    \item 输出：10 个类别的分数向量
\end{itemize}

\begin{figure}[H]
\centering
\includegraphics[width=0.95\textwidth]{slides-images/slide-005.jpg}
\caption{线性分类器的三种理解方式：代数视角、模板视角和几何视角\protect\footnotemark}
\end{figure}
\footnotetext{Slides 第6页。}

\begin{knowledgebox}{线性分类器的三种视角}
\begin{itemize}
    \item \textbf{代数视角}：矩阵乘法 $W\mathbf{x} + \mathbf{b}$，每行独立计算一个类别的分数
    \item \textbf{模板视角}：将每行权重还原为图像形状，可视化每个类别学到的``模板''
    \item \textbf{几何视角}：每行权重定义一条决策边界（$W_i \cdot \mathbf{x} + b_i = 0$），将输入空间划分为不同区域
\end{itemize}
几何视角的优势在于可以直观判断线性模型的能力边界——如果数据无法被一条直线分开，线性模型就无能为力。
\end{knowledgebox}

\subsection{损失函数}

有了模型和分数，如何衡量模型的好坏？这就需要损失函数。给定数据集 $\{(x_i, y_i)\}_{i=1}^N$，总损失定义为：

$$L(W) = \frac{1}{N}\sum_{i=1}^{N} L_i(f(x_i, W), y_i)$$

\begin{figure}[H]
\centering
\includegraphics[width=0.95\textwidth]{slides-images/slide-010.jpg}
\caption{损失函数：Data Loss 衡量模型预测与训练数据的匹配程度\protect\footnotemark}
\end{figure}
\footnotetext{Slides 第14页。}

\textbf{Softmax 损失}（交叉熵损失）是最常用的分类损失函数：

$$L_i = -\log\left(\frac{e^{s_{y_i}}}{\sum_j e^{s_j}}\right)$$

\begin{itemize}
    \item 先对分数做指数变换使其全为正
    \item 再除以总和得到概率分布（所有值之和为 1）
    \item 取正确类别概率的负对数作为损失
\end{itemize}

\begin{importantbox}{Softmax 的作用}
Softmax 将任意一组浮点数转换为概率分布：所有输出值之和为 1，且每个值表示对应类别的概率。分数越高的类别获得越大的概率值。当正确类别的概率接近 1 时损失接近 0；当正确类别的概率很低时损失很大。
\end{importantbox}

\subsection{本章小结}

图像分类是计算机视觉的核心任务，面临语义鸿沟、光照变化等多重挑战。线性分类器通过 $f = W\mathbf{x} + \mathbf{b}$ 将图像映射到类别分数，Softmax 损失函数衡量预测质量。接下来的问题是：如何在损失函数中加入正则化以防止过拟合，以及如何找到最优的 $W$。

%% ========== 正则化 ==========

